\documentclass[11pt,a4paper]{article}
\usepackage{../common/td_style}

\begin{document}

\makeuniversityheader{Corrigé des Travaux Dirigés : Série 05}{Analyse en Composantes Principales (ACP) en Biologie}{\textcolor{BioTeal}{\textbf{CORRIGÉ DÉTAILLÉ}}}

\begin{solutionbox}{Exercice 1 : Centrage-Réduction \& Nécessité de la Standardisation}
\begin{enumerate}[label=\textbf{\alph*.}]
  \item \textbf{Formule du Z-score :}
  \begin{equation*}
    z_{ij} = \frac{x_{ij} - \bar{x}_j}{s_j}
  \end{equation*}
  Cette transformation soustrait la moyenne (centrage) et divise par l'écart-type (réduction), donnant une variable adimensionnelle de moyenne 0 et de variance 1.

  \item \textbf{Calcul des coordonnées centrées-réduites du patient :}
  \begin{itemize}
    \item Glycémie : $z_{i1} = \frac{1.60 - 1.10}{0.25} = \frac{0.50}{0.25} = \mathbf{+2.00}$.
    \item Transaminases ALAT : $z_{i2} = \frac{105.0 - 45.0}{30.0} = \frac{60.0}{30.0} = \mathbf{+2.00}$.
    \item Cholestérol : $z_{i3} = \frac{6.80 - 5.20}{0.80} = \frac{1.60}{0.80} = \mathbf{+2.00}$.
  \end{itemize}
  \textit{Interprétation :} Le patient se situe exactement à $+2$ écarts-types de la moyenne pour les 3 paramètres. En termes d'anomalie relative standardisée, les trois déviations sont rigoureusement équivalentes (environ 97.7\textsuperscript{e} percentile de la population normale).

  \item \textbf{Effet catastrophique d'une ACP non normée :}
  La variance brute de l'ALAT est de $s_2^2 = 900\,\text{UI}^2/\text{L}^2$, alors que celle de la glycémie n'est que de $s_1^2 = 0.0625\,\text{g}^2/\text{L}^2$.
  \par Le rapport de variance est de $\frac{900}{0.0625} = 14\,400$ !
  L'ACP non normée maximiserait la variance brute : le premier axe factoriel serait confondu à $99.9\%$ avec l'ALAT simplement à cause du choix arbitraire de l'unité de mesure (UI/L vs g/L). La glycémie et le cholestérol seraient totalement écrasés et ignorés.

  \item \textbf{Conclusion biologique :}
  En biologie médicale et en métabolomique omique, les variables mesurent des grandeurs physiques aux échelles incomparables. L'**ACP normée (sur matrice de corrélation)** est impérative car elle attribue à chaque biomarqueur une variance égale à 1, garantissant un poids analytique équilibré.
\end{enumerate}
\end{solutionbox}

\begin{solutionbox}{Exercice 2 : Matrice de Corrélation, Valeurs Propres \& Éboulis de Kaiser}
\begin{enumerate}[label=\textbf{\alph*.}]
  \item \textbf{Somme totale des valeurs propres :}
  \begin{equation*}
    \sum_{j=1}^5 \lambda_j = 2.65 + 1.25 + 0.60 + 0.35 + 0.15 = \mathbf{5.00}
  \end{equation*}
  Dans une ACP normée de dimension $p$, la matrice de corrélation a des 1 sur sa diagonale principale. La trace de $R$ vaut : $\text{Tr}(R) = \sum r_{jj} = p \times 1 = p = 5$. Comme la somme des valeurs propres est égale à la trace de la matrice, on retrouve exactement $\sum \lambda_j = 5$.

  \item \textbf{Proportions de variance expliquée et pourcentages cumulés :}
  \begin{center}
    \small
    \begin{tabular}{lccc}
      \toprule
      \textbf{Composante (Axe)} & \textbf{Valeur propre ($\lambda_k$)} & \textbf{\% Inertie ($\lambda_k / 5 \times 100$)} & \textbf{\% Cumulé} \\
      \midrule
      \textbf{PC1 ($F_1$)} & 2.65 & $\frac{2.65}{5} \times 100 = \mathbf{53.0\%}$ & $53.0\%$ \\
      \textbf{PC2 ($F_2$)} & 1.25 & $\frac{1.25}{5} \times 100 = \mathbf{25.0\%}$ & $\mathbf{78.0\%}$ \\
      \textbf{PC3 ($F_3$)} & 0.60 & $\frac{0.60}{5} \times 100 = 12.0\%$ & $90.0\%$ \\
      \textbf{PC4 ($F_4$)} & 0.35 & $\frac{0.35}{5} \times 100 = 7.0\%$  & $97.0\%$ \\
      \textbf{PC5 ($F_5$)} & 0.15 & $\frac{0.15}{5} \times 100 = 3.0\%$  & $100.0\%$ \\
      \bottomrule
    \end{tabular}
  \end{center}

  \item \textbf{Application du critère de Kaiser ($\lambda \ge 1.0$) :}
  On ne retient que les axes dont la valeur propre est supérieure ou égale à 1 (ceux qui expliquent plus de variance qu'une variable individuelle standardisée).
  Ici, seuls $\lambda_1 = 2.65$ et $\lambda_2 = 1.25$ sont $\ge 1$. On retient donc exactement **$k = 2$ axes factoriels**.

  \item \textbf{Inertie du premier plan factoriel (PC1 -- PC2) :}
  Le plan factoriel (PC1, PC2) capture $53.0\% + 25.0\% = \mathbf{78.0\%}$ de l'information totale. Réduire 5 dimensions à un plan 2D préserve l'essentiel de la structure biologique tout en éliminant le bruit expérimental résiduel ($22\%$).
\end{enumerate}
\end{solutionbox}

\begin{solutionbox}{Exercice 3 : Interprétation du Cercle des Corrélations \& Métriques}
\begin{enumerate}[label=\textbf{\alph*.}]
  \item \textbf{Tableau complété des qualités de représentation ($\cos^2$) :}
  \begin{center}
    \resizebox{\linewidth}{!}{%
    \begin{tabular}{lccccc}
      \toprule
      \textbf{Biomarqueur} & \textbf{PC1} & \textbf{PC2} & \textbf{$\cos^2(\text{PC1})$} & \textbf{$\cos^2(\text{PC2})$} & \textbf{Qualité Globale ($\cos^2_{1+2}$)} \\
      \midrule
      \textbf{MDA}          & $+0.88$ & $-0.20$ & $0.88^2 = \mathbf{0.774}$ & $(-0.20)^2 = \mathbf{0.040}$ & $0.774 + 0.040 = \mathbf{0.814}$ \\
      \textbf{IL-6}         & $+0.92$ & $-0.15$ & $0.92^2 = \mathbf{0.846}$ & $(-0.15)^2 = \mathbf{0.023}$ & $0.846 + 0.023 = \mathbf{0.869}$ \\
      \textbf{TNF-$\alpha$} & $+0.85$ & $-0.25$ & $0.85^2 = \mathbf{0.723}$ & $(-0.25)^2 = \mathbf{0.063}$ & $0.723 + 0.063 = \mathbf{0.786}$ \\
      \textbf{Catalase}     & $-0.80$ & $-0.45$ & $(-0.80)^2 = \mathbf{0.640}$ & $(-0.45)^2 = \mathbf{0.203}$ & $0.640 + 0.203 = \mathbf{0.843}$ \\
      \textbf{SOD}          & $-0.10$ & $+0.85$ & $(-0.10)^2 = \mathbf{0.010}$ & $0.85^2 = \mathbf{0.723}$    & $0.010 + 0.723 = \mathbf{0.733}$ \\
      \bottomrule
    \end{tabular}%
    }
  \end{center}
  Toutes les variables ont un $\cos^2_{1+2} > 0.70$ : elles sont toutes très bien projetées dans ce plan factoriel.

  \item \textbf{Variables directrices des axes :}
  \begin{itemize}
    \item **Axe PC1 :** Directement tiré par l'**IL-6** ($\cos^2 = 0.846$) et le **MDA** ($\cos^2 = 0.774$) d'un côté, et par la **Catalase** ($\cos^2 = 0.640$) de l'autre.
    \item **Axe PC2 :** Presque exclusivement dominé par la **Superoxyde Dismutase (SOD)** ($\cos^2 = 0.723$).
  \end{itemize}

  \item \textbf{Relation IL-6 et MDA (Angle très aigu) :}
  L'angle $\theta \approx 0^\circ$ signifie une corrélation linéaire positive très forte ($r \approx \cos(0^\circ) = +1$). Biologiquement, l'inflammation cytokinique (IL-6) et la peroxydation lipidique membranaire (MDA) évoluent de concert : l'orage cytokinique stimule la production d'espèces réactives de l'oxygène (ROS) qui dégradent les lipides polyinsaturés.

  \item \textbf{Opposition MDA/IL-6 vs Catalase (Angle de $180^\circ$) \& Dénomination :}
  La Catalase est en opposition diamétrale ($r \approx -1$) avec les biomarqueurs de stress oxydatif : plus le tissu s'enflamme et peroxyle, plus ses défenses enzymatiques antioxydantes s'épuisent.
  \par **Dénomination biologique de PC1 :** « \textbf{Axe du Stress Oxydatif \& de l'Épuisement Antioxydant} ».
\end{enumerate}
\end{solutionbox}

\begin{solutionbox}{Exercice 4 : Plan des Individus \& Discrimination Clinique}
\begin{enumerate}[label=\textbf{\alph*.}]
  \item \textbf{Nature non supervisée de l'ACP :}
  L'ACP n'a jamais reçu les étiquettes de diagnostic (Témoins / Diabète / NASH) lors du calcul : elle a simplement cherché les directions de variance maximale dans le nuage de points.
  \par Le fait que les trois groupes cliniques se séparent spontanément et de manière ordonnée le long de PC1 (Témoins à gauche $\to$ Diabète au centre $\to$ NASH à droite) démontre que les altérations métaboliques mesurées reflètent fidèlement la progression physiopathologique de la maladie.

  \item \textbf{Profil du nouveau patient $(\text{PC1} = +2.40, \, \text{PC2} = -0.30)$ :}
  Avec un score $\text{PC1} = +2.40$ très élevé, ce patient se situe dans la région extrême des patients atteints de **Stéato-Hépatite Non Alcoolique (NASH sévère)** avec stress oxydatif massif et effondrement des défenses antioxydantes.
  \par *Examens de confirmation recommandés :* Élastographie hépatique (FibroScan) pour quantifier la fibrose, bilan hépatique complet (ALAT/ASAT, $\gamma$-GT) et biopsie hépatique de contrôle.

  \item \textbf{Différenciation par l'axe PC2 :}
  Pour un même niveau d'atteinte globale ($\text{PC1} = +1.80$), le patient à $\text{PC2} = +1.50$ possède une activité **SOD très élevée** (réponse adaptative de protection encore active), tandis que le patient à $\text{PC2} = -1.20$ souffre d'un épuisement enzymatique complet en superoxyde dismutase (pronostic hépatique plus sévère).
\end{enumerate}
\end{solutionbox}

\begin{solutionbox}{Exercice 5 : Sélection de Biomarqueurs \& Élimination de la Redondance}
\begin{enumerate}[label=\textbf{\alph*.}]
  \item \textbf{Redondance analytique entre IL-6 et TNF-$\alpha$ :}
  Les deux vecteurs sont quasi-superposés sur le cercle des corrélations (angle $< 5^\circ$), ce qui indique une corrélation colinéaire proche de $r = 0.95$. Doser les deux revient à payer deux fois le coût d'anticorps ELISA pour mesurer exactement la même information biologique sous-jacente.

  \item \textbf{Choix entre IL-6 et TNF-$\alpha$ :}
  On choisit sans hésiter l'**IL-6**, car sa qualité de représentation sur l'axe principal est supérieure : $\cos^2(\text{PC1})_{\text{IL-6}} = 0.846$ contre $0.723$ pour le TNF-$\alpha$. L'IL-6 est plus informative et plus discriminante pour le phénotypage des patients.

  \item \textbf{Panel minimal optimal (2 ou 3 biomarqueurs) :}
  Pour capturer l'ensemble des 2 axes indépendants ($78\%$ de la variance) :
  \begin{enumerate}
    \item **IL-6** (ou MDA) : pour capturer l'axe inflammatoire majeur PC1.
    \item **Catalase** : pour mesurer la capacité de détoxification antioxydante (pôle opposé de PC1).
    \item **SOD** : indispensable car c'est la seule variable qui capture l'axe orthogonal PC2.
  \end{enumerate}
  Ce panel à 3 marqueurs permet de reconstituer l'intégralité du profil factoriel du patient à une fraction du coût d'un profilage complet.
\end{enumerate}
\end{solutionbox}

\end{document}
